\section{Architecture Identity}

\begin{itemize}
  \item Architecture name: VIV-32.
  \item Toolchain and implementation prefix: \texttt{viv32}.
  \item Architecture style: 32-bit load/store architecture.
  \item Primary implementation targets:
  \begin{itemize}
    \item reference emulator;
    \item QEMU target;
    \item SystemVerilog/Verilog hardware implementation.
  \end{itemize}
\end{itemize}

\section{Word and Address Model}

\begin{itemize}
  \item Word size: 32 bits.
  \item Address size: 32 bits.
  \item Byte size: 8 bits.
  \item Endianness: big-endian.
  \item Memory model: flat byte-addressed memory.
  \item Address space model: flat 32-bit physical address space.
  \item Virtual memory / paging: not supported in VIV-32 v\version{}.
  \item Segmentation: not supported in VIV-32 v\version{}.
\end{itemize}

\section{Instruction Model}

\begin{itemize}
  \item Instruction width: fixed 32-bit instructions.
  \item Instruction encoding style: regular RISC-style encoding with a small
        number of instruction format families.
  \item Variable-width instructions: not supported in VIV-32 v\version{}.
  \item Compressed or variable-width instruction encodings are deferred to a
        later architecture extension.
  \item Separate register tables for general purpose and floating point
        operations.
  \item Load/store architecture: yes.
  \item Arithmetic and logical operations operate on registers.
  \item Memory is accessed only through explicit load and store instructions.
  \item Memory operand policy: arithmetic and logical instructions do not
        operate directly on memory operands.
  \item Immediate sizes: defined per instruction format.
  \item Branch offset unit: 32-bit instruction words, shifted left by two bits
        to form byte offsets.
  \item Branch range: BF-type branches use \code{offset22}; BC-type
        branches use \code{offset14}.
  \item Opcode extension policy: reserve opcode space for future architectural
        extensions.
  \item Possible future extensions include:
  \begin{itemize}
    \item compressed instructions;
    \item atomic instructions;
    \item privileged memory-management features;
    \item SIMD, vector, matrix, or co-processor operations.
  \end{itemize}
\end{itemize}

\section{Register Model}

\begin{itemize}
  \item Sixteen 32-bit general-purpose registers with signed and unsigned
        operations.
  \item Eight 64-bit floating-point registers.
  \item General-purpose registers are named \code{\$0} through \code{\$15}.
  \item Floating-point registers are named \code{\#0} through
        \code{\#7}.
  \item Register \code{\$0} is hardwired to zero; reads always return zero, and
        writes are ignored.
  \item Register \code{\$13} is reserved by ABI convention as an optional frame
        pointer \code{\$fp}.
  \item Register \code{\$14} is reserved by ABI convention as the stack pointer
        \code{\$sp}.
  \item Register \code{\$15} is used as the link register \code{\$lr} for call
        and return instructions.
  \item The frame pointer is not a dedicated architectural register.
  \item The program counter is directly readable using a dedicated instruction.
  \item The program counter is not directly writable as a general-purpose
        register.
  \item Control flow changes are performed through:
  \begin{itemize}
    \item branch instructions;
    \item jump instructions;
    \item call instructions;
    \item return instructions;
    \item trap instructions;
    \item interrupt-return instructions.
  \end{itemize}
  \item Special/control registers include:
        \begin{itemize}
        \item \code{\%pc}: program counter;
        \item \code{\%sr}: status register;
        \item \code{\%epc}: exception program counter;
        \item \code{\%esr}: exception status register, containing the
              status-register value saved on exception entry;
        \item \code{\%ecause}: exception cause;
        \item \code{\%edata}: exception-specific data associated with the most
              recent exception or interrupt.
        \item \code{\%evbase}: exception vector base.
        \end{itemize}
  \item Parallel/vector register support is not included in the base ISA.
  \item Future extensions may define SIMD, vector, matrix, or co-processor
        facilities.
  \item Matrix or co-processor support is deferred to a future extension or
        memory-mapped accelerator design.
\end{itemize}

\section{Arithmetic and Flags}

\begin{itemize}
\item Integer representation: two's complement.
\item Floating-point register representation: IEEE-754 binary64.
\item Floating-point memory representation: IEEE-754 binary32 and binary64.
\item Signed and unsigned integer arithmetic results wrap modulo \(2^{32}\).
\item Floating-point arithmetic is performed using IEEE-754 binary64 precision.
\item Floating-point arithmetic and floating-point conversions use
      round-to-nearest, ties-to-even unless otherwise specified by an
      instruction.
\item VIV-32 defines condition flags in the processor status register.
\item Condition flags are \code{NZCVEIUX}:
      \begin{itemize}
      \item \code{NEGATIVE} (\code{N}): set when the numeric result of a
            flag-setting operation is negative. For floating-point results,
            the sign of zero does not cause \code{N} to be set.
      \item \code{ZERO} (\code{Z}): set when the numeric result of a
            flag-setting operation is zero. For floating-point operations,
            both \code{+0.0} and \code{-0.0} are considered zero.
      \item \code{CARRY} (\code{C}): used by integer arithmetic. Set when an
            unsigned addition produces a carry out of bit 31. For subtraction,
            \code{C} is set when no borrow is required and cleared when a borrow
            is required. Floating-point operations clear \code{C} unless otherwise
            specified.
      \item \code{OVERFLOW} (\code{V}): used by integer arithmetic. Set when
            the mathematical result cannot be represented in the destination
            integer domain. Floating-point operations clear \code{V} unless
            otherwise specified.
      \item \code{ARITHMETIC ERROR} (\code{E}): set when an arithmetic or
            conversion operation encounters an exceptional or invalid
            condition that does not trap. This includes integer
            divide-by-zero, floating-point divide-by-zero, invalid
            floating-point operations, unordered flag-setting floating-point
            comparisons, and invalid numeric conversions.
      \item \code{INFINITY} (\code{I}): used by floating-point operations.
            Set when the resulting IEEE-754 value is positive or negative
            infinity. Integer operations clear \code{I} unless otherwise
            specified.
      \item \code{UNDERFLOW} (\code{U}): used by floating-point operations.
            Set when a floating-point result is too small in magnitude to be
            represented normally in the destination floating-point format.
            Integer operations clear \code{U} unless otherwise specified.
      \item \code{INEXACT} (\code{X}): used by floating-point operations and
            conversions. Set when the mathematically exact result cannot be
            represented exactly in the destination format and rounding is
            therefore required. Integer operations clear \code{X} unless
            otherwise specified.
      \end{itemize}
\item Integer arithmetic operations update the relevant \code{NZCVE} flags
      and clear \code{IUX} unless otherwise specified.
\item Floating-point arithmetic operations update the relevant
      \code{NZEIUX} flags and clear \code{CV} unless otherwise specified.
\item Floating-point divide-by-zero sets \code{E}; if the resulting value is
      infinity, \code{I} is also set.
\item Invalid floating-point operations producing \code{NaN} set \code{E}.
\item Flag-setting floating-point comparisons involving \code{NaN} set
      \code{E} to indicate an unordered comparison.
\item Floating-point to integer conversion of \code{NaN}, infinity, or a
      value outside the destination integer range sets \code{E}.
\item Floating-point narrowing from binary64 to binary32 updates
      \code{NZIUX} according to the binary32 result and sets \code{E} only
      when the conversion is otherwise invalid.
\item Integer comparisons update the relevant integer comparison flags and
      clear floating-point-specific flags unless otherwise specified.
\item Flag-setting floating-point comparisons update the relevant
      floating-point flags; \code{C} and \code{V} have no floating-point
      comparison meaning.
\item Flag-based branch instructions test the current processor-status flags.
      They do not interpret the flags as signed, unsigned, or floating-point
      comparisons and do not perform a comparison themselves.
\end{itemize}

\section{Memory Access}

\begin{itemize}
  \item Supported load sizes:
  \begin{itemize}
    \item 8-bit byte;
    \item 16-bit halfword;
    \item 32-bit word;
  \end{itemize}
  \item Supported store sizes:
  \begin{itemize}
    \item 8-bit byte;
    \item 16-bit halfword;
    \item 32-bit word;
  \end{itemize}
  \item Byte accesses may occur at any address.
  \item Halfword accesses must be aligned to a 2-byte boundary.
  \item Word accesses must be aligned to a 4-byte boundary.
  \item Misaligned access behaviour: misaligned halfword or word accesses raise
        a\\
        \code{MisalignedAccessException}.
  \item VIV-32 v\version{} defines no cache architecture.
  \item All memory accesses are treated as direct accesses to memory or
        memory-mapped devices.
  \item VIV-32 uses memory-mapped I/O.
  \item Device registers occupy reserved physical address ranges defined by the
        platform specification.
\end{itemize}

\section{Control Flow}

\begin{itemize}
  \item Jump mechanism: \PC-relative immediate jumps (\code{jmp}) and
        register-indirect jumps (\code{jr}).
  \item Call mechanism: \PC-relative call with link register \code{call}, plus
        register-indirect call (\code{jalr}).
  \item Link register behaviour: \code{call} instructions store the address of
        the next instruction in \code{\%lr}. \code{jalr} provides the option to
        use a different return register.
  \item Return mechanism: \code{ret} jumps to the address held in the link
        register \code{\%lr}. A general register-indirect jump instruction may
        also be used for returns, if previously loaded with an address via
        \code{jalr}.
  \item \code{trap} and \code{syscall} instructions transfer control to their
        associated exception vectors and record the continuation address and cause in
        exception control registers \code{\%epc} and \code{\%ecause}.
        \code{trap} additionally stores its immediate operand in
        \code{\%edata}. \code{syscall} sets \code{\%edata} to \code{0}.
  \item Halt instruction behaviour: the \code{halt} instruction stops
        instruction execution until reset or platform-defined external action.
  \item NOP encoding: \code{0x00000000}.
  \item Extended immediates: VIV-32 v\version{} does not use implicit
        next-word operands in the base instruction stream. Full-width constants
        and absolute addresses can be constructed using the pseudo-instructions
        \code{li} and \code{la}, respectively.
\end{itemize}

\section{Reset, Exceptions and Interrupts}

\begin{itemize}
  \item Privilege model: VIV-32 v\version{} defines machine mode only. User and
        supervisor modes are reserved for future versions.
  \item Reset: reset begins directly at \code{0x00}.
  \item Exception vector model: VIV-32 uses a vector base register
        \code{\%evbase} and cause-specific vector offsets. The exact vector
        layout is defined later in this document.
  \item Interrupt vector model: interrupts transfer control through the
        exception vector mechanism.
  \item Exception types:

        \begin{itemize}
          \item \code{Reset};
          \item \code{IllegalInstructionException};
          \item \code{MisalignedFetchException};
          \item \code{MisalignedAccessException};
          \item \code{SoftwareTrap};
          \item \code{SystemCall};
          \item \code{Interrupt};
          \item \code{BusError}.
        \end{itemize}
        
  \item Interrupt enable/disable mechanism: the processor status register
        contains an interrupt-enable bit, which can be cleared (disabled) with
        \code{di} and set (enabled) with \code{ei}. The CPU starts with this
        bit cleared---a program wishing to use interrupts should enable this
        after configuring the relevant peripheral settings.
  \item \code{iret} instruction: VIV-32 defines an interrupt-return instruction
        which resumes execution from the saved exception program counter in
        \code{\%epc}.
  \item Illegal instruction behaviour: illegal instructions raise an
        \code{IllegalInstructionException}.
  \item External interrupt: reserved for platform devices.
  \item Timer interrupt: supported in VIV-32 v\version{} through the general
        platform device interrupt ISR.
  \item Divide-by-zero behaviour: divide-by-zero does not trap in VIV-32
        v\version{}. Software is responsible for checking \code{ARITHMETIC ERROR}.
\end{itemize}
